% 原分章PDF106页附录7.G标题—110页末，印刷页682—686。
\section*{附录7.G\quad 带噪声线性模型的迹方法}
\addcontentsline{toc}{section}{附录7.G 带噪声线性模型的迹方法}
\markright{附录7.G 带噪声线性模型的迹方法}

回顾
\begin{equation}
\Sigma_{XY}=\Sigma_XA^T\tag{7.G.1}
\end{equation}
以及
\begin{equation}
\Sigma_Y=A\Sigma_XA^T+\lambda I,\tag{7.G.2}
\end{equation}
其中 $A$ 为正交矩阵。由式（7.112），有
\begin{equation}
\begin{aligned}
\widetilde A&=\Sigma_XA^T(A\Sigma_XA^T+\lambda I)^{-1}\\
 &=\Sigma_X(\Sigma_X+\lambda I)^{-1}A^T.
\end{aligned}\tag{7.G.3}
\end{equation}
假设协方差矩阵 $\Sigma_X$ 的特征值分解为
\begin{equation}
\Sigma_X=U\Lambda U^T,\tag{7.G.4}
\end{equation}
其中 $U^TU=I$、$\Lambda=\operatorname{diag}(\mu_1,\ldots,\mu_n)$，且 $\mu_1\ge\mu_2\ge\cdots\mu_n\ge0$。于是，由式（7.G.3）、（7.G.4），得到
\begin{equation}
\begin{aligned}
\widetilde A&=U\Delta U^T(U\Delta U^T+\lambda I)^{-1}A^T\\
 &=U\Delta(\Delta+\lambda I)^{-1}U^TA^T.
\end{aligned}\tag{7.G.5}
\end{equation}
因此，
\begin{equation}
\begin{aligned}
\tau_n(\widetilde A\widetilde A^T)
 &=\tau_n\bigl(U\Delta(\Delta+\lambda I)^{-1}U^TA^TAU
  (\Delta+\lambda I)^{-1}\Delta U^T\bigr)\\
 &=\tau_n\left(\bigl[\Delta(\Delta+\lambda I)^{-1}\bigr]^2\right).
\end{aligned}\tag{7.G.6}
\end{equation}
利用式（7.G.2）、（7.G.4），得到
\begin{equation}
\Sigma_Y=AU\Delta U^TA^T+\lambda I,\tag{7.G.7}
\end{equation}
从而
\begin{equation}
\tau_n(\Sigma_Y)=\tau_n(\Delta)+\lambda.\tag{7.G.8}
\end{equation}
现在考虑模型
\[
X=\widetilde AY+\widetilde e,
\]
并计算 $\tau_n(\widetilde A\Sigma_Y\widetilde A^T)$：
\begin{equation}
\begin{aligned}
&\tau_n(\widetilde A\Sigma_Y\widetilde A^T)\\
&\quad=\tau_n\Bigl(U\Delta(\Delta+\lambda I)^{-1}U^TA^T
 (AU\Delta U^TA^T+\lambda I)\\
&\hspace{8em}{}\cdot AU(\Delta+\lambda I)^{-1}\Delta U^T\Bigr)\\
&\quad=\tau_n\bigl(U\Delta(\Delta+\lambda I)^{-1}
 (\Delta+\lambda I)(\Delta+\lambda I)^{-1}\Delta U^T\bigr)\\
&\quad=\tau_n\bigl(U\Delta(\Delta+\lambda I)^{-1}\Delta U^T\bigr)\\
&\quad=\tau_n\bigl(\Delta^2(\Delta+\lambda I)^{-1}\bigr).
\end{aligned}\tag{7.G.9}
\end{equation}
由式（7.G.6）、（7.G.8）、（7.G.9），得到评估因果方向 $Y\to X$ 的统计量：
\begin{equation}
\begin{aligned}
\Delta_{Y\to X}
={}&\log\frac{\tau_n(\widetilde A\Sigma_Y\widetilde A^T)}
 {\tau_n(\Sigma_Y)\tau_n(\widetilde A\widetilde A^T)}\\
={}&\log\frac{\tau_n\bigl(\Delta^2(\Delta+\lambda I)^{-1}\bigr)}
 {(\tau_n(\Delta)+\lambda)\tau_n\left(\bigl[\Delta(\Delta+\lambda I)^{-1}\bigr]^2\right)}\\
={}&\log\frac1n\sum_{i=1}^n\frac{\mu_i^2}{\mu_i+\lambda}
 -\log\frac1n\sum_{i=1}^n(\mu_i+\lambda)\\
&{}-\log\frac1n\sum_{i=1}^n\frac{\mu_i^2}{(\mu_i+\lambda)^2}.
\end{aligned}\tag{7.G.10}
\end{equation}
现在用归纳法证明 $\Delta_{Y\to X}>0$。令 $n=2$，需证明
\begin{equation}
\frac12\left(\frac{\mu_1^2}{\mu_1+\lambda}+\frac{\mu_2^2}{\mu_2+\lambda}\right)
 >\frac12(\mu_1+\lambda+\mu_2+\lambda)\frac12
 \left(\frac{\mu_1^2}{(\mu_1+\lambda)^2}+\frac{\mu_2^2}{(\mu_2+\lambda)^2}\right).\tag{7.G.11}
\end{equation}
事实上，式（7.G.11）蕴含
\begin{equation}
2\left(\frac{\mu_1^2}{\mu_1+\lambda}+\frac{\mu_2^2}{\mu_2+\lambda}\right)
 >(\mu_1+\lambda+\mu_2+\lambda)
 \left(\frac{\mu_1^2}{(\mu_1+\lambda)^2}+\frac{\mu_2^2}{(\mu_2+\lambda)^2}\right).\tag{7.G.12}
\end{equation}
展开式（7.G.12）的右侧，得到
\begin{equation}
\begin{aligned}
&(\mu_1+\lambda+\mu_2+\lambda)
 \left(\frac{\mu_1^2}{(\mu_1+\lambda)^2}+\frac{\mu_2^2}{(\mu_2+\lambda)^2}\right)\\
&\quad=\frac{\mu_1^2}{\mu_1+\lambda}+\frac{\mu_2^2}{\mu_2+\lambda}
 +(\mu_1+\lambda)\frac{\mu_2^2}{(\mu_2+\lambda)^2}
 +(\mu_2+\lambda)\frac{\mu_1^2}{(\mu_1+\lambda)^2}.
\end{aligned}\tag{7.G.13}
\end{equation}
结合式（7.G.12）、（7.G.13），需要证明
\begin{equation}
\left(\frac{\mu_1^2}{\mu_1+\lambda}+\frac{\mu_2^2}{\mu_2+\lambda}\right)
 >(\mu_1+\lambda)\frac{\mu_2^2}{(\mu_2+\lambda)^2}
 +(\mu_2+\lambda)\frac{\mu_1^2}{(\mu_1+\lambda)^2}.\tag{7.G.14}
\end{equation}
注意
\begin{equation}
\begin{aligned}
\frac{\mu_1^2}{\mu_1+\lambda}-(\mu_2+\lambda)\frac{\mu_1^2}{(\mu_1+\lambda)^2}
 &=\left(1-\frac{\mu_2+\lambda}{\mu_1+\lambda}\right)\frac{\mu_1^2}{\mu_1+\lambda}\\
 &=\frac{\mu_1-\mu_2}{\mu_1+\lambda}\frac{\mu_1^2}{\mu_1+\lambda}
\end{aligned}\tag{7.G.15}
\end{equation}
以及
\begin{equation}
\begin{aligned}
\frac{\mu_2^2}{\mu_2+\lambda}-(\mu_1+\lambda)\frac{\mu_2^2}{(\mu_2+\lambda)^2}
 &=\left(1-\frac{\mu_1+\lambda}{\mu_2+\lambda}\right)\frac{\mu_2^2}{\mu_2+\lambda}\\
 &=\frac{\mu_2-\mu_1}{\mu_2+\lambda}\frac{\mu_2^2}{\mu_2+\lambda}.
\end{aligned}\tag{7.G.16}
\end{equation}
由式（7.G.15）、（7.G.16），得到
\begin{equation}
\begin{aligned}
&\left(\frac{\mu_1^2}{\mu_1+\lambda}+\frac{\mu_2^2}{\mu_2+\lambda}\right)
 -(\mu_1+\lambda)\frac{\mu_2^2}{(\mu_2+\lambda)^2}
 -(\mu_2+\lambda)\frac{\mu_1^2}{(\mu_1+\lambda)^2}\\
&\quad=\frac{\mu_1-\mu_2}{\mu_1+\lambda}\frac{\mu_1^2}{\mu_1+\lambda}
 +\frac{\mu_2-\mu_1}{\mu_2+\lambda}\frac{\mu_2^2}{\mu_2+\lambda}\\
&\quad=(\mu_1-\mu_2)\left(\frac{\mu_1^2}{(\mu_1+\lambda)^2}
 -\frac{\mu_2^2}{(\mu_2+\lambda)^2}\right)\\
&\quad=(\mu_1-\mu_2)\left(\frac{\mu_1}{\mu_1+\lambda}-\frac{\mu_2}{\mu_2+\lambda}\right)
 \left(\frac{\mu_1}{\mu_1+\lambda}+\frac{\mu_2}{\mu_2+\lambda}\right)\\
&\quad=\frac{\lambda(\mu_1-\mu_2)^2}{(\mu_1+\lambda)(\mu_2+\lambda)}
 \left(\frac{\mu_1}{\mu_1+\lambda}+\frac{\mu_2}{\mu_2+\lambda}\right)>0.
\end{aligned}\tag{7.G.17}
\end{equation}
这就证明了不等式（7.G.11）。现在假设 $n=k$ 时，下列不等式成立：
\begin{equation}
\frac1n\sum_{i=1}^n\frac{\mu_i^2}{\mu_i+\lambda}
 >\left(\frac1n\sum_{i-1}^n(\mu_i+\lambda)\right)
 \left(\frac1n\sum_{i=1}^n\frac{\mu_i^2}{(\mu_i+\lambda)^2}\right).\tag{7.G.18}
\end{equation}
下面证明 $n=k+1$ 时，不等式（7.G.18）仍然成立。令
\[
S_k^1=\sum_{i=1}^k\frac{\mu_i^2}{\mu_i+\lambda},\qquad
S_k^2=\sum_{i=1}^k(\mu_i+\lambda),\qquad
S_k^3=\sum_{i=1}^k\frac{\mu_i^2}{(\mu_i+\lambda)^2}.
\]
于是，不等式（7.G.18）表明
\begin{equation}
kS_k^1>S_k^2S_k^3\tag{7.G.19}
\end{equation}
成立。注意
\[
\begin{aligned}
\sum_{i=1}^{k+1}\frac{\mu_i^2}{\mu_i+\lambda}
 &=S_k^1+\frac{\mu_{k+1}^2}{\mu_{k+1}+\lambda},\\
\sum_{i=1}^{k+1}(\mu_i+\lambda)&=S_k^2+\mu_{k+1}+\lambda,\qquad\text{且}\\
\sum_{i=1}^{k+1}\frac{\mu_i^2}{(\mu_i+\lambda)^2}
 &=S_K^3+\frac{\mu_{k+1}^2}{(\mu_{k+1}+\lambda)^2}.
\end{aligned}
\]
为了说明不等式（7.G.18）成立，必须证明
\[
\frac1{k+1}\sum_{i=1}^{k+1}\frac{\mu_i^2}{\mu_i+\lambda}
 >\left(\frac1{k+1}\sum_{i=1}^{k+1}(\mu_i+\lambda)\right)
 \left(\frac1{k+1}\sum_{i=1}^{k+1}\frac{\mu_i^2}{(\mu_i+\lambda)^2}\right),
\]
即
\begin{equation}
(k+1)\sum_{i=1}^{k+1}\frac{\mu_i^2}{\mu_i+\lambda}
 >\left(\sum_{i=1}^{k+1}(\mu_i+\lambda)\right)
 \left(\sum_{i=1}^{k+1}\frac{\mu_i^2}{(\mu_i+\lambda)^2}\right).\tag{7.G.20}
\end{equation}
注意
\begin{equation}
\sum_{i=1}^{k+1}\frac{\mu_i^2}{\mu_i+\lambda}
 =S_k^1+\frac{\mu_{k+1}^2}{\mu_{k+1}+\lambda},\tag{7.G.21}
\end{equation}
\begin{equation}
\sum_{i=1}^{k+1}(\mu_i+\lambda)=S_k^2+\mu_{k+1}+\lambda,\tag{7.G.22}
\end{equation}
\begin{equation}
\sum_{i=1}^{k+1}\frac{\mu_i^2}{(\mu_i+\lambda)^2}
 =S_k^3+\frac{\mu_{k+1}^2}{(\mu_{k+1}+\lambda)^2}.\tag{7.G.23}
\end{equation}
将式（7.G.21）—（7.G.23）代入式（7.G.20），得到
\begin{equation}
(k+1)S_k^1+(k+1)\frac{\mu_{k+1}^2}{\mu_{k+1}+\lambda}
 >(S_k^2+\mu_{k+1}+\lambda)
 \left(S_k^3+\frac{\mu_{k+1}^2}{(\mu_{k+1}+\lambda)^2}\right).\tag{7.G.24}
\end{equation}
利用式（7.G.19），可将式（7.G.24）化为
\begin{equation}
S_k^1+k\frac{\mu_{k+1}^2}{\mu_{k+1}+\lambda}
 >(\mu_{k+1}+\lambda)S_k^3
 +\frac{\mu_{k+1}^2}{(\mu_{k+1}+\lambda)^2}S_k^2.\tag{7.G.25}
\end{equation}
将式（7.G.25）的右侧移至左侧，得到
\begin{equation}
\begin{aligned}
&\sum_{i=1}^k\frac{\mu_i^2}{\mu_i+\lambda}
 \left(1-\frac{\mu_{k+1}+\lambda}{\mu_i+\lambda}\right)
 +\sum_{i=1}^k\frac{\mu_{k+1}^2}{\mu_{k+1}+\lambda}
 \left(1-\frac{\mu_i+\lambda}{\mu_{k+1}+\lambda}\right)\\
&\quad=\sum_{i=1}^k\frac{\mu_i^2}{\mu_i+\lambda}
 \frac{\mu_i-\mu_{k+1}}{\mu_i+\lambda}
 +\sum_{i=1}^k\frac{\mu_{k+1}^2}{\mu_{k+1}+\lambda}
 \frac{\mu_{k+1}-\mu_i}{\mu_{k+1}+\lambda}\\
&\quad=\sum_{i=1}^k(\mu_i-\mu_{k+1})
 \left(\frac{\mu_i^2}{(\mu_i+\lambda)^2}
 -\frac{\mu_{k+1}^2}{(\mu_{k+1}+\lambda)^2}\right)\\
&\quad=\sum_{i=1}^k(\mu_i-\mu_{k+1})
 \left(\frac{\mu_i}{\mu_i+\lambda}-\frac{\mu_{k+1}}{\mu_{k+1}+\lambda}\right)
 \left(\frac{\mu_i}{\mu_i+\lambda}+\frac{\mu_{k+1}}{\mu_{k+1}+\lambda}\right)\\
&\quad=\sum_{i=1}^k(\mu_i-\mu_{k+1})^2\lambda
 \frac1{(\mu_{i+1}+\lambda)(\mu_{k+1}+\lambda)}
 \left(\frac{\mu_i}{\mu_i+\lambda}+\frac{\mu_{k+1}}{\mu_{k+1}+\lambda}\right)\ge0.
\end{aligned}\tag{7.G.26}
\end{equation}
结合式（7.G.20）、（7.G.24）—（7.G.26），证明了
\begin{equation}
\frac1n\sum_{i=1}^n\frac{\mu_i^2}{\mu_i+\lambda}
 >\left(\frac1n\sum_{i=1}^n(\mu_i+\lambda)\right)
 \left(\frac1n\sum_{i=1}^n\frac{\mu_i^2}{(\mu_i+\lambda)^2}\right)\tag{7.G.27}
\end{equation}
以及
\begin{equation}
\begin{aligned}
\Delta_{Y\to X}
={}&\log\frac1n\sum_{i=1}^n\frac{\mu_i^2}{\mu_i+\lambda}
 -\log\frac1n\sum_{i=1}^n(\mu_i+\lambda)\\
&{}-\log\frac1n\sum_{i=1}^n\frac{\mu_i^2}{(\mu_i+\lambda)^2}>0.
\end{aligned}\tag{7.G.28}
\end{equation}
